\section{Exact attainment of source weights}
\label{subsec:valuation-extremal-weights}

We prove Theorem~\ref{V:invariant-main} by constructing the small input
orders required by Theorem~\ref{V:exact}. Let $X$ be smooth connected
projective with smoothly semipositive $L=-K_X$, and assume
$H^i(X,\mathcal O_X)=0$ for all $i>0$. Let an algebraic torus $T$
act effectively on $X$. If $X$ is a point, its constant section proves
the theorem. Otherwise use $G_0$ and $\ell$ as in
\eqref{V:null-functional}. The anticanonical line $L=\det T_X$ has its
natural linearization.  The action on sections is the left action
induced by the action on the total space.

\subsection{Source weights and cohomology}

Fix a cocharacter $\lambda:\mathbb G_m\to T$.  Its unique source
component $S$ attracts a dense open subset as $t\to0$, by the
Bia\l ynicki-Birula decomposition \cite{BB1973}; see also the smooth attracting decomposition in \cite[Theorem~1.5 and Corollary~7.3]{JelisiejewSienkiewicz2019}.  Let $c$ be
the weight of $L$ on $S$.  This is the sum of the positive normal
tangent weights, so $c>0$ for a nonzero cocharacter.  For the zero
cocharacter take $S=X$ and $c=0$.  The rational attracting map to $S$
is dominant.  Resolve it and pull back holomorphic forms; birational
invariance and Hodge symmetry give
$H^i(S,\mathcal O_S)=0$ for $i>0$.

Average the logarithmic weights of the metric over the circle in
$\lambda(\mathbb G_m)$.  For a nonzero vector $u\in L_x$, put
$q_u(s)=\log|\lambda(e^s)u|$.  The logarithmic norm of this
nonvanishing holomorphic orbit section is superharmonic by semipositive
curvature.  Circle invariance therefore makes $q_u$ concave in $s$.
There is a smooth function $\theta:X\to\mathbb R$ such that
$q_u'(s)=\theta(\lambda(e^s)x)$; it is independent of the scaling
of $u$ and equals the fiber weight on each fixed component.  If $x$
is in the dense attracting set, continuity gives
$q_u'(s)\to c$ as $s\to-\infty$.  Concavity then gives
$\theta(x)\le c$.  Density and continuity extend this inequality
to all of $X$, so every fixed-component fiber weight is at most $c$.

For all sufficiently large integers $l$, the weight $(l-1)c$ occurs in
some $H^q(X,(l-1)L)$.  Indeed, expand the equivariant holomorphic
Lefschetz formula at infinity in its weight variable $z$.
For a normal tangent weight $w$, the inverse normal factor is
$(1-z^{-w}e^{-x})^{-1}$.  If $w>0$, it has constant term one and lower
powers; if $w<0$, all its powers are strictly negative.  Every component
other than the source has a negative normal weight.  Maximality of $c$
therefore makes the coefficient of $z^{(l-1)c}$ precisely
$\chi(S,(l-1)L|_S)$.  This polynomial in $l-1$ has value
$\chi(S,\mathcal O_S)=1$ at zero, and is nonzero for all large $l$.
See \cite{AtiyahSegal1968,AtiyahSinger1968} for the fixed-component formula.

Equivariant hard Lefschetz with the smooth semipositive coefficient
$lL$ \cite[Theorem~0.1]{DPS2001} gives a surjection
\[
 H^0(X,\Omega_X^{n-q}\otimes lL)\longrightarrow
 H^q(X,K_X+lL)=H^q(X,(l-1)L).
\]
A circle-invariant K\"ahler form makes this map equivariant.  On an
infinite subsequence we obtain nonzero eigensections $u_l$ of
$\Omega_X^a\otimes lL$, with one fixed $a$, of weight $(l-1)c$ for
$\lambda$.  Decompose each into $T$-weight components and select a
nonzero component; all the components have the same $\lambda$-weight.
Thus the $u_l$ can be chosen to be $T$-eigensections.

\subsection{Determinants with bounded orbit order}

The use of saturated determinants follows the method of
Lazi\'c--Peternell \cite[Lemma~4.1]{LP2018} and its anticanonical
adaptation in \cite[Lemma~5.1]{LMPTX2023}. Those statements include
a numerical-nontriviality hypothesis on the twisting line. Here the
positive unbounded twists are already available, and the generic-span
and wedge construction below proves exactly the assertion needed,
also for a numerically trivial $L$.

Each section $u_l$ defines a map $L^{-l}\to\Omega_X^a$. Take the
span of their image lines over $\C(X)$, of rank $e>0$, and saturate
its determinant line in $\bigwedge^e\Omega_X^a$. This construction is
independent of a choice of rational frame for $L$.  Denote the resulting
line bundle by $M$.  Rank-one saturation in a locally free sheaf is
reflexive and hence a line bundle on the smooth variety; its natural
linearization extends from the subbundle locus across codimension two.
The needed curvature-intersection sign is $\ell(M)\leq0$; it is
automatic when $\nu=n$.  Suppose $\nu<n$.  The smooth metric on
$L=-K_X$ determines a positive volume form whose Ricci form is
$\alpha$.  For every $b>0$, scale this volume to the volume of the
K\"ahler class $G_0+bL$ and apply the Calabi--Yau theorem
\cite{Yau1978}.  The resulting K\"ahler metric has Ricci form
$\alpha\ge0$.  K\"ahler curvature symmetries identify the mean
curvature of its tangent bundle with the Ricci endomorphism.  Duals
and exterior powers therefore make the mean curvature of
$\bigwedge^e\Omega_X^a$ nonpositive.

Apply the subbundle curvature inequality to the inclusion of $M$
where it is nonzero.  Its logarithmic norm extends as a locally
quasiplurisubharmonic function across the zero set.  Its Hessian has
a smooth lower bound there, so the singular trace measure is
nonnegative and the extension preserves the degree inequality.  Hence
\[
 M\cdot(G_0+bL)^{n-1}\leq0.
\]
Since $L^{\nu+1}=0$, all powers above $b^\nu$ vanish in this
intersection polynomial.  The coefficient of $b^\nu$ is a positive
multiple of $\ell(M)$, so letting $b\to\infty$ proves the required
sign.

Choose a fixed finite collection of the $u_l$ whose image lines form
a generic basis. Every further nonzero $u_l$ can be completed to a basis
by $e-1$ members of this collection. There are finitely many choices;
on an infinite subsequence one choice is fixed. Taking these wedges
gives nonzero eigensections of $M+m_jL$, where $m_j$ is the sum of the
$e$ twists and tends to infinity. Their weights are $(m_j-e)c$, since
each of the $e$ constituent forms contributes its twist minus one.  Their zero divisors $N_j$
satisfy $\ell(N_j)=0$, by effectivity and $\ell(M)\leq0$.
Let $v$ be the order obtained by substituting $\lambda(t)x$ for the
generic point $x$ of $X$.  Properness extends this orbit to a formal arc
at zero, with center at the generic point of $S$; its order is a
discrete valuation up to rescaling, or is zero.  If $c_M$ is the fiber
weight of $M$ at $S$, then
\[
 v(N_j)=m_jc+c_M-(m_j-e)c=c_M+ec.
\]
This is Lemma~\ref{V:order-weight} for the declared left action. The
constant $c_M+ec$ retains the determinant-rank shift.

\subsection{Exact source weights and an invariant product}

Theorem~\ref{V:exact} now gives an effective rational
representative of $L$ of order zero at the generic source.  After
clearing denominators, some section of $mL$ restricts nontrivially to
$S$.  Its weight decomposition consequently contains a normalized
weight whose value on $\lambda$ is exactly $c$.  If $T$ is trivial,
this is already a nonzero invariant section.  Assume from now on that
$T$ has positive dimension.

The function $c$ is the maximum of the finitely many $T$-fixed fiber
weights evaluated on $\lambda$: the source is $T$-stable and contains
a $T$-fixed point by the Borel fixed-point theorem.  It thus extends
to a continuous rational polyhedral homogeneous function on real
cocharacters.  It is strictly positive away from zero.  Indeed a
nontrivial zero or negative locus would, by rational polyhedrality,
contain a nonzero rational cocharacter, contrary to the source-weight
calculation.  Every normalized section weight is bounded above by $c$
on each cocharacter, by the same orbit-order test, and the preceding
argument attains this upper bound.  The comparison between fixed-fiber
and section-weight polytopes for semiample linearized line bundles,
and the resulting invariant-product criterion, are established in
M\"uller~\cite[Proposition~3.13 and Corollary~3.14]{Muller2025}.
Here exact order attainment supplies the supporting section weights.

Testing a lattice basis of cocharacters and its negatives bounds
all normalized section weights.  Their closed convex hull is therefore
compact; its support function equals $c$ on rational cocharacters by
exact attainment, and then on all real cocharacters by continuity.
Strict positivity of $c$ away from zero puts zero in its interior.
Choose a small simplex surrounding zero in that convex hull and
approximate its vertices by finite convex combinations of actual
normalized weights.  Sufficiently close approximations still surround
zero, so the finitely many actual weights in these combinations have
zero in their convex hull.  Lemma~\ref{V:convex-product} now multiplies finitely many corresponding
eigensections to produce a nonzero invariant section.

This proves Theorem~\ref{V:invariant-main}. For an arbitrary torus action, replace $T$ by its effective quotient
$T/\ker(T\to\operatorname{Aut}(X))$. The entire kernel, including
its finite part, acts trivially on the natural anticanonical line, so
a section invariant under the quotient is invariant under $T$.
